\documentclass[10pt]{article}
\providecommand{\FIGDIR}{fig}
\newcommand{\DIGESTDATE}{14 August 2026}
% ===================== PM Research Watch - shared preamble =====================
\usepackage[T1]{fontenc}
\usepackage[utf8]{inputenc}
\usepackage{amsmath,amssymb}
\usepackage{tgpagella}
\usepackage{tgheros}
\renewcommand{\sfdefault}{qhv}
\renewcommand{\ttdefault}{lmtt}
\usepackage{microtype}
\usepackage[english]{babel}

\usepackage[a4paper,top=17mm,bottom=17mm,left=17mm,right=17mm,
            headsep=6mm,footskip=10mm]{geometry}
\usepackage{graphicx}
\usepackage{xcolor}
\usepackage{colortbl}
\usepackage{tikz}
\usetikzlibrary{calc}
\usepackage[most]{tcolorbox}
\usepackage{titlesec}
\usepackage{fancyhdr}
\usepackage{booktabs}
\usepackage{array}
\usepackage{longtable}
\usepackage{enumitem}
\usepackage{multicol}
\usepackage{ragged2e}
\usepackage[hidelinks]{hyperref}

% ------------------------------------------------------------------ palette
\definecolor{Ink}    {HTML}{14262E}
\definecolor{Deep}   {HTML}{0F4C5C}
\definecolor{Teal}   {HTML}{1F7A8C}
\definecolor{Sky}    {HTML}{4FA3C4}
\definecolor{Sage}   {HTML}{6FA287}
\definecolor{Amber}  {HTML}{E3A008}
\definecolor{Clay}   {HTML}{B5643C}
\definecolor{Coral}  {HTML}{D1495B}
\definecolor{Violet} {HTML}{6A5B9E}
\definecolor{Slate}  {HTML}{4A6572}
\definecolor{Paper}  {HTML}{FBF9F5}
\definecolor{Mist}   {HTML}{EEEAE2}
\definecolor{Rule}   {HTML}{D9D3C9}

\pagecolor{Paper}
\color{Ink}

\hypersetup{colorlinks=true, linkcolor=Deep, urlcolor=Teal, citecolor=Teal,
            pdftitle={PM Research Watch}, pdfauthor={Automated daily digest}}

% ------------------------------------------------------------------ headings
\titleformat{\section}
  {\sffamily\bfseries\large\color{Deep}}{}{0pt}
  {}[{\vspace{1.2mm}\color{Amber}\titlerule[1.1pt]}]
\titlespacing*{\section}{0pt}{5.5mm}{2.8mm}

\titleformat{\subsection}
  {\sffamily\bfseries\normalsize\color{Teal}}{}{0pt}{}
\titlespacing*{\subsection}{0pt}{3.6mm}{1.4mm}

\setlist[itemize]{leftmargin=4.2mm, itemsep=1.1mm, topsep=1.3mm, parsep=0pt,
                  label=\textcolor{Teal}{\footnotesize$\blacktriangleright$}}

% ------------------------------------------------------------------ header/footer
\pagestyle{fancy}
\fancyhf{}
\renewcommand{\headrulewidth}{0pt}
\renewcommand{\footrulewidth}{0pt}
\fancyhead[L]{\sffamily\scriptsize\color{Slate}PM RESEARCH WATCH\;\textcolor{Rule}{\textbar}\;\DIGESTDATE}
\fancyhead[R]{\sffamily\scriptsize\color{Slate}Particulate matter monitoring \& health effects}
\fancyfoot[C]{\sffamily\scriptsize\color{Slate}\thepage}

% ------------------------------------------------------------------ components
% metric tile
\newcommand{\metric}[3]{%
  \begin{tcolorbox}[enhanced, colback=white, colframe=Rule, boxrule=0.5pt,
    arc=1.6mm, left=2.2mm, right=2.2mm, top=1.8mm, bottom=1.8mm,
    borderline west={1.7mm}{0pt}{#3}, nobeforeafter, valign=center,
    height=17mm, width=\linewidth]
    \RaggedRight
    {\sffamily\bfseries\fontsize{17}{17}\selectfont\color{#3}#1}\\[0.5mm]
    {\sffamily\fontsize{6.5}{7.4}\selectfont\color{Slate}#2}
  \end{tcolorbox}}

% pull-quote / key insight box
\newtcolorbox{keybox}[1][]{enhanced, colback=Mist, colframe=Mist,
  arc=1.6mm, left=3.2mm, right=3.2mm, top=2.4mm, bottom=2.4mm,
  borderline west={1.7mm}{0pt}{Amber}, #1}

% subtopic band
\newcommand{\band}[2]{%
  \vspace{3.2mm}\noindent
  \begin{tcolorbox}[enhanced, colback=#2, colframe=#2, arc=1.2mm,
    left=2.6mm, right=2.6mm, top=1.1mm, bottom=1.1mm, boxrule=0pt,
    width=\linewidth, nobeforeafter]
    {\sffamily\bfseries\small\color{white}#1}
  \end{tcolorbox}\vspace{1.6mm}}

% one paper entry: {short cite}{title}{journal}{doi}{body}
\newcommand{\paperentry}[5]{%
  \par\vspace{1.9mm}\noindent
  {\sffamily\bfseries\footnotesize\color{Ink}#1}\;%
  {\sffamily\fontsize{7}{8}\selectfont\color{Slate}#3}\\[0.4mm]
  {\footnotesize\itshape\color{Deep}#2}\\[0.6mm]
  {\footnotesize #4}\\[0.3mm]
  {\fontsize{6.6}{7.6}\selectfont\color{Slate}\texttt{doi:}\,\href{https://doi.org/#5}{\texttt{\detokenize{#5}}}}
  \par}

% caption style
\newcommand{\figcap}[1]{\par\vspace{1.0mm}
  {\fontsize{7.2}{8.6}\selectfont\color{Slate}\textsf{\textbf{Fig.}}\; #1\par}\vspace{1.5mm}}

\providecommand{\permil}{\textperthousand}

\setlength{\parindent}{0pt}
\setlength{\parskip}{0pt}
\renewcommand{\arraystretch}{1.18}


\begin{document}

% ==================================================================== MASTHEAD
\thispagestyle{empty}
\begin{tikzpicture}[remember picture, overlay]
  \fill[Deep] (current page.north west) rectangle
      ($(current page.north east)+(0,-67mm)$);
  \fill[Amber] ($(current page.north west)+(0,-67mm)$) rectangle
      ($(current page.north east)+(0,-68.4mm)$);
  \foreach \x/\y/\r in {22/12/0.9, 41/26/1.5, 63/9/0.7, 88/31/1.1, 112/17/1.9,
                        138/29/0.8, 159/13/1.3, 178/54/1.0, 196/21/1.6,
                        30/44/0.6, 74/57/1.2, 125/48/0.9, 168/6/0.8, 191/59/0.7,
                        52/18/1.0, 100/60/0.6, 148/20/1.1, 66/49/1.4, 205/45/0.8}
    \fill[Sky, opacity=0.30] ($(current page.north west)+(\x mm,-\y mm)$) circle (\r mm);
\end{tikzpicture}

\vspace*{4mm}
{\sffamily\fontsize{7.6}{9}\selectfont\color{Sky}%
 AUTOMATED DAILY LITERATURE DIGEST\;\textbullet\;PhD RESEARCH WATCH\par}
\vspace{2.2mm}
{\sffamily\bfseries\fontsize{27}{29}\selectfont\color{white}%
 Particulate Matter\par}
{\sffamily\fontsize{27}{29}\selectfont\color{Sky}%
 Monitoring \& Health Effects\par}
\vspace{5.0mm}
{\sffamily\fontsize{8.4}{10}\selectfont\color{white}%
 Issue 2026-08-14\;\;\textcolor{Amber}{\textbar}\;\;Entry window 14 August 2026\;\;%
 \textcolor{Amber}{\textbar}\;\;14 records carried, 17 rejected\par}
\vspace{18mm}

% -------------------------------------------------------------- metric strip
\noindent
\begin{minipage}[t]{0.190\linewidth}\metric{14}{RECORDS CARRIED\\(17 REJECTED)}{Deep}\end{minipage}\hfill
\begin{minipage}[t]{0.190\linewidth}\metric{9}{SUBTOPIC CLUSTERS\\REPRESENTED}{Teal}\end{minipage}\hfill
\begin{minipage}[t]{0.190\linewidth}\metric{11}{EFFECT ESTIMATES\\WITH INTERVALS}{Sky}\end{minipage}\hfill
\begin{minipage}[t]{0.190\linewidth}\metric{5}{INSTRUMENTATION /\\EXPOSURE RECORDS}{Amber}\end{minipage}\hfill
\begin{minipage}[t]{0.190\linewidth}\metric{0}{METADATA-ONLY\\(NO ABSTRACT)}{Clay}\end{minipage}

\vspace{6mm}

% -------------------------------------------------------------- signal of day
\section*{\sffamily\bfseries\color{Deep}Signal of the day}
\vspace{-1.5mm}{\color{Amber}\rule{\linewidth}{1.1pt}}\vspace{2.5mm}

\begin{keybox}
{\sffamily\bfseries\small\color{Deep}Two records landed on the same day that, read
together, put a number on the exposure-misclassification ceiling --- and it is low
enough to matter.}\\[1.4mm]
\footnotesize
Makacha et al.\ (\textit{Front Glob Womens Health}) ran 5-day-per-season minute-resolved
\emph{personal} PM$_{2.5}$ monitoring on 343 parous women in The Gambia, Kenya and
Mozambique --- 3{,}190 person-days --- against reference monitors at nearby health
facilities. Ambient-to-personal concordance was \textbf{0.01--0.56 with regression slopes
below 1 in every setting}, and exposure was consistently higher \emph{away} from home,
peaking around 12:00--16:00. A fixed-site monitor is not a weak proxy for personal dose in
these settings; across the worst site pairings it is barely a proxy at all.

Set that beside Atalell et al.\ (\textit{Environ Res}), who detect developmental
vulnerability at school entry in 54{,}458 Western Australian children at a mean PM$_{2.5}$
of \textbf{7.40--7.54\,$\mu$g/m$^3$} --- below the WHO annual guideline --- with
\textbf{OR 1.06 (1.03--1.08) for vulnerability in two or more domains per
1\,$\mu$g/m$^3$}. The effect size per unit is small, so the inference rests entirely on
exposure contrast being measured cleanly. One study says effects are detectable at a
1\,$\mu$g/m$^3$ resolution; the other says the assignment method used to get there can be
off by more than that. The gap between them is the case for dense, personal-scale,
time-resolved measurement, and it is a measurement argument rather than an epidemiological
one.
\end{keybox}

\vspace{3mm}

% -------------------------------------------------------------- exec brief
\section*{\sffamily\bfseries\color{Deep}What changed today}
\vspace{-1.5mm}{\color{Amber}\rule{\linewidth}{1.1pt}}\vspace{2.5mm}

\footnotesize
\begin{itemize}
\item \textbf{A dependence-aware meta-analysis erases a cancer signal.} Negahi et al.\
pooled 25 estimates from 14 studies of climate-sensitive exposures and cancer. The
conventional random-effects estimate was \textbf{1.07 (1.02--1.12)}; under robust
variance estimation, accounting for multiple estimates drawn from shared cohorts, it
became \textbf{1.02 (0.93--1.11)}. $I^2 = 97\%$, GRADE low or very low in all ten groups.
This is a methodological result more than a substantive one, and it should be cited
whenever a pooled air-pollution estimate is built from overlapping cohorts.

\item \textbf{Downscaling needs bias correction \emph{before} super-resolution.}
d'Hondt et al.\ (preprint) argue that existing deep-learning downscalers implicitly treat
the coarse CTM field as a blurred copy of the fine field, conflating cross-resolution
\emph{model} differences with genuinely missing spatial detail. Separating structural
correction from spatial refinement is directly relevant to any LCS-plus-CTM fusion
scheme, where the sensor network supplies exactly the detail the coarse model lacks.

\item \textbf{A data-quality index that scores a low-cost sensor at 58/100.} Bajaj \&
Settle decompose fitness-for-use into accuracy, calibration, completeness,
representational consistency and timeliness. On five July days of OpenAQ Los Angeles data
a reference monitor scored $\approx$99 and an LCS $\approx$58 --- penalised not mainly for
accuracy but for completeness, outliers during a 4 July episode, and reporting that
stopped altogether. The framing is worth borrowing even though the demonstration is
n$\,$=$\,$2.

\item \textbf{Composition, not mass, again carries the toxicology.} Ayala et al.\ resolved
54 Catalan PM$_{10}$ filter samples into four MCR-ALS components and regressed six
SH-SY5Y endpoints against them ($R^2_{CV}$ 0.53--0.75). Winter biomass-burning markers
(levoglucosan, dehydroabietic acid, PAHs) drove ROS, p53 activation and cell death;
biogenic SOA tracers were unimportant, and a spring--summer SOA component inhibited
acetylcholinesterase \emph{without} cytotoxicity. Two distinct neurotoxic pathways from
one mass metric.

\item \textbf{An interventional record worth tracking rather than reading.} Wittkopp et
al.\ published the protocol (NCT05994937) for a double-blind sham-controlled bedroom HEPA
trial in adults with prediabetes, with continuous glucose monitoring as the primary
outcome and continuous indoor PM$_{2.5}$ throughout. No results yet; logged to
\texttt{trials.json} for the weekly watch.
\end{itemize}

\clearpage

% ==================================================================== ANALYTICS
\section*{\sffamily\bfseries\color{Deep}Corpus analytics}
\vspace{-1.5mm}{\color{Amber}\rule{\linewidth}{1.1pt}}\vspace{3mm}

\noindent\includegraphics[width=\linewidth]{\FIGDIR/f1_subtopics.png}
\figcap{Subtopic assignment of the 14 carried records, single-label by dominant
contribution. Nine of the ten clusters are occupied --- the broadest spread in a week ---
but no cluster reaches five. \textit{Sensing, forecasting \& instrumentation} leads at
4; \textit{Neuro / mental health} is the only empty cluster, which is unusual and is an
artefact of the day rather than a trend.}

\vspace{2mm}
\noindent\includegraphics[width=\linewidth]{\FIGDIR/f2_design_endpoint.png}
\figcap{Left --- study architecture. Observational designs (4), modelling (3) and
measurement campaigns (3) split most of the corpus; the two records in the
\textit{Trial / intervention} column are a protocol and a feasibility pilot, so neither
reports an outcome. Right --- health endpoint. Respiratory endpoints lead at 3, and six
records carry no human health endpoint --- the four sensing records, the personal-exposure
campaign, and the second-hand-smoke pilot, whose reported outcome is indoor PM$_{2.5}$
rather than a clinical measure.}

\vspace{2mm}
\noindent\includegraphics[width=\linewidth]{\FIGDIR/f3_metric_geography.png}
\figcap{Left --- particle metric. PM$_{2.5}$ dominates, but two records depart from mass
entirely: particle number size distribution and CCN in the Lhasa campaign, and a burden
estimate carrying no PM metric at all (CO/NO$_2$/SO$_2$ only). Right --- geographic
provenance. China contributes four records and Europe three; sub-Saharan Africa is
represented by a first-party personal-monitoring campaign rather than, as is more usual
in this watch, a review of its absence.}

\vspace{2mm}
\noindent\includegraphics[width=\linewidth]{\FIGDIR/f6_lifecourse.png}
\figcap{Life-course windows addressed. Counts are non-exclusive and describe this issue
only. The in-utero window is occupied for the first time since 10 August (Shanghai cord
blood; personal monitoring in pregnancy), and puberty/adolescence is empty --- the
mirror image of the 13 August distribution.}

\vspace{2mm}
\noindent\includegraphics[width=\linewidth]{\FIGDIR/f4_heatmap.png}
\figcap{Subtopic $\times$ study-architecture cross-tabulation. The matrix is almost
perfectly diagonal --- nine subtopics spread across seven architecture columns with no
column holding more than three records --- which is the visual signature of a broad,
shallow day. The instrumentation $\leftrightarrow$ epidemiology silo flagged on 12 and 13
August persists: every \textit{Sensing} and \textit{Exposure} record sits in a
measurement or modelling column and none carries a health endpoint. Ayala et al.\ are the
partial exception and the only record today in which exposure and outcome were measured
by the same group on the same samples.}

\vspace{2mm}
\noindent\includegraphics[width=\linewidth]{\FIGDIR/f5_forest.png}
\figcap{The 11 interval estimates carried today, on a common logarithmic axis.
\textbf{The exposure contrasts are not comparable across sources.} Atalell et al.\ report
per-1\,$\mu$g/m$^3$ ORs from a whole-of-population cohort, which is why their intervals
are the tightest on the plot and their point estimates the closest to unity; Chen et al.\
report per-unit adjusted ORs from 101 mother--infant pairs, which is why theirs are the
widest; Zhai \& Meng report per-IQR relative risks; and the three Negahi et al.\ estimates
are pooled meta-analytic risk ratios, two of which are fitted to the same underlying
studies under different dependence assumptions and are therefore not independent of one
another. Makrufardi et al.\ report six odds ratios without intervals and are not plotted.}

\clearpage

% ==================================================================== DIGEST
\section*{\sffamily\bfseries\color{Deep}Paper-level digest}
\vspace{-1.5mm}{\color{Amber}\rule{\linewidth}{1.1pt}}\vspace{1mm}

{\fontsize{7.4}{9}\selectfont\color{Slate}Ordered by cluster. Each entry gives the
methodological core and the finding that matters, not an abstract paraphrase.\par}

\band{Sensing, forecasting \& instrumentation\hfill 4 records}{Amber}

\paperentry{d'Hondt et al.\ 2026}
{Neural downscaling of air-quality simulations requires structural correction before spatial refinement}
{Research Square (preprint)}
{Deep-learning downscaling of CAMS-class chemical transport model output, which is
capped near 10\,km for continental forecasts. The argument is that existing approaches
either treat the coarse field as a smoothed copy of the fine field or learn a direct
coarse-to-fine map, and both conflate two different things: cross-resolution differences
in \emph{model behaviour}, and genuinely missing spatial detail. Separating a structural
correction step from the super-resolution step is what the paper proposes and tests.
Directly transferable to LCS-plus-CTM fusion, where the network supplies the detail the
model lacks and the model supplies the physics the network lacks --- but a bias-corrected
coarse field is the precondition. Preprint, not peer reviewed; carried on relevance,
and no numeric claim from it is used in this issue's analytics.}
{10.21203/rs.3.rs-10627504/v1}

\paperentry{Bajaj \& Settle 2026}
{Rapid assessment of air pollution data quality using a composite index}
{Environ Monit Assess}
{A Data Quality Index aggregating five 0--100 dimensions --- accuracy (30\%), calibration
(20\%), completeness (20\%), representational consistency (15\%), timeliness (15\%) ---
with weights the authors state should vary by stakeholder. Demonstrated on hourly OpenAQ
PM records for Los Angeles, 1--5 July 2025: a reference-grade monitor scored
\textbf{$\approx$99} and a low-cost sensor \textbf{$\approx$58}, the latter penalised for
calibration bias, low completeness, extreme outliers during the 4 July episode, and zero
timeliness once reporting ceased on 5 July. The value is in keeping the components
legible rather than collapsing to a single opaque score. The demonstration is
n$\,$=$\,$2 over five days and includes a fireworks episode --- an outlier regime chosen,
apparently, for illustration --- so the numbers are a worked example, not a benchmark.}
{10.1007/s10661-026-15780-6}

\paperentry{Yaqoob et al.\ 2026}
{Semiparametric hybrid air quality forecasting using Prophet ensemble and feature selection}
{Sci Rep}
{Prophet models the trend and seasonality of daily PM$_{2.5}$; the residuals are then
modelled by bagging, boosting, stacking and voting regressors, with Lasso, SHAP and
recursive feature elimination applied to control the high-dimensional lag/interaction
feature space. Tested on Seoul daily data and a high-dimensional simulated dataset
against XGBoost, SVR and LSTM baselines; Prophet + boosting with Lasso and SHAP gave the
lowest RMSE. The residual-modelling decomposition is a sensible way to keep an
interpretable trend component, but the abstract reports no RMSE values, no forecast
horizon and no comparison against a persistence baseline, which is the benchmark that
actually discriminates in daily PM forecasting.}
{10.1038/s41598-026-64590-w}

\paperentry{Zhang et al.\ 2026}
{High frequency of urban new particle formation on the Tibetan Plateau}
{Atmos Chem Phys}
{In-situ particle number size distribution measurements at Tibet University, Lhasa
(3{,}650\,m), across spring and autumn campaigns, with days classified NPF / Undefined /
Non-Event from PNSD evolution plus size-segregated counts and growth rates. NPF event
frequency was \textbf{86.2\% in autumn and 62.1\% in spring} --- extremely high by urban
standards --- and estimated CCN was $(2.3\pm1.5)\times10^3$\,cm$^{-3}$ on NPF days versus
$(1.3\pm0.7)\times10^3$ on non-event days. Autumn events were associated with greater
sulfur-precursor availability; no dominant control was identifiable in spring, which the
authors state plainly. Relevant here as a reminder that number-based instrumentation
resolves a regime mass-based monitoring cannot see at all; the high-altitude,
low-condensation-sink setting limits transfer to lowland urban networks.}
{10.5194/acp-26-11473-2026}

\band{Exposure assessment \& modelling\hfill 1 record}{Amber}

\paperentry{Makacha et al.\ 2026}
{Spatio-temporal and structural determinants of personal PM$_{2.5}$ exposure in three sub-Saharan African geographies}
{Front Glob Womens Health}
{Five days per season of minute-resolved personal PM$_{2.5}$ monitoring on 343 parous
women in The Gambia, Kenya and Mozambique, 3{,}190 person-days, aligned to
microenvironment, fuel use, occupation, wealth and meteorology, with facility-based
ambient monitors for comparison. Daily medians were modest --- 10.4 (7.6--14.7),
8.1 (6.8--9.8) and 10.9 (9.3--13.5)\,$\mu$g/m$^3$ --- but episodic spikes dominated, and
WHO 24-h exceedance was strongly seasonal (The Gambia \textbf{55.9\% dry vs 13.6\% wet}).
\textbf{Ambient--personal concordance was 0.01--0.56 with slopes $<$1 throughout.}
Exposure was higher away from home with a 12:00--16:00 peak, so clean-fuel transition
alone will not close the gap. The 5-day-per-season design is thin for characterising
episodic exposure and the authors do not report how much of the between-person variance
it captures.}
{10.3389/fgwh.2026.1843221}

\band{Reproductive \& developmental\hfill 2 records}{Sky}

\paperentry{Atalell et al.\ 2026}
{Low-level PM$_{2.5}$ and developmental vulnerability at school entry, Western Australia}
{Environ Res}
{Whole-of-population linkage of the Australian Early Development Census (2009, 2012, 2015
waves) to midwives' notifications, birth registrations and hospital morbidity data;
54{,}458 children assigned daily PM$_{2.5}$ at Statistical Area Level 2 across prenatal,
early-childhood and cumulative windows; Bayesian multilevel logistic regression against
the 2009 10th-percentile domain cut-offs. Mean exposure \textbf{7.40--7.54\,$\mu$g/m$^3$},
below WHO guidance. Early-childhood exposure: \textbf{OR 1.04 (1.01--1.07)} for
vulnerability in one or more domains and \textbf{1.06 (1.03--1.08)} in two or more, per
1\,$\mu$g/m$^3$, with signals in physical health, social competence and language/cognitive
domains. Prenatal and cumulative windows were null. The authors concede the critical
limitation themselves: prenatal and early-childhood exposures are highly correlated, so
the window attribution is not identified. Area-level assignment at SA2 resolution over a
1\,$\mu$g/m$^3$ contrast is the second limitation, and the one this watch cares about.}
{10.1016/j.envres.2026.125493}

\paperentry{Chen et al.\ 2026}
{Cord blood DNA methylation as a mediator between prenatal air pollution and infant allergic symptoms}
{Pediatr Allergy Immunol}
{Shanghai prospective birth cohort, 101 mother--infant pairs with spatiotemporally
modelled residential exposure, caregiver-reported cumulative allergic symptoms to
6 months, and targeted bisulfite sequencing of na\"ive CD4$^+$ T cells in a subset of
\textbf{30} cord-blood samples. First-trimester PM (\textbf{aOR 2.27, 1.09--4.98}) and
NO$_2$ (\textbf{2.11, 1.02--4.53}) were associated with allergic symptoms; third-trimester
sulfate, nitrate and ammonium components were positively associated; a CpG in the
\textit{PRKACB} promoter showed statistical evidence of mediation for the sulfate
association, supported by functional assays. The mediation arm rests on 30 samples and
the outcome is caregiver-reported at 6 months, which is the weakest available allergic
phenotype. The authors label the findings exploratory, correctly. Carried for the
component-resolved exposure design, not the effect sizes.}
{10.1111/pai.70424}

\band{Respiratory \& allergic\hfill 2 records}{Sage}

\paperentry{Makrufardi et al.\ 2026}
{Ambient relative humidity and particulate matter vs.\ asthma treatment step}
{World Allergy Organ J}
{288 newly diagnosed adults with asthma followed at least 6 months under GINA, with
treatment step --- not symptoms or spirometry --- as the outcome, and both daily means and
short-term \emph{differences} in RH, temperature and PM$_{2.5}$ as exposures. A
1\,$\mu$g/m$^3$ increase in the 7-day PM$_{2.5}$ difference was associated with 1.052-fold
odds of step 4 treatment and 1.056-fold odds of medium-dose ICS + LABA; a
1\,$\mu$g/m$^3$ increase in the 30-day difference gave 1.205-fold odds of triple therapy
or biologics. Associations were clearer in the cold season. \textbf{No confidence
intervals are reported in the abstract}, which is why none of these appear in the forest
plot. Using variability rather than level as the exposure is the interesting choice ---
it maps onto what a high-time-resolution sensor measures well and a monthly average does
not --- but with a single centre and 288 patients this is explicitly
hypothesis-generating, as the authors say.}
{10.1016/j.waojou.2026.101444}

\paperentry{Zhai \& Meng 2026}
{Apparent temperature, air pollutants and pneumonia hospitalisation in rural Western China}
{Int J Environ Health Res}
{Distributed lag non-linear models on NRCMS surveillance data from Jiuquan, 2015--2020,
with stratification and explicit pollutant $\times$ temperature interaction terms. Apparent
temperature showed an inverted-U with a peak \textbf{RR 1.50 (1.41--1.60)}. Per-IQR
increases gave \textbf{CO 1.11 (1.01--1.21)} and \textbf{PM$_{10}$ 1.05 (1.00--1.12)} as
the only intervals excluding or touching unity; PM$_{2.5}$ \textbf{1.04 (0.97--1.12)},
SO$_2$, O$_3$ and NO$_2$ were null. Synergy was reported between low apparent temperature
and high pollution, and among pollutants. Risks were higher in children under 7 and in
autumn/winter. A rural setting is a genuine gap-filler --- most Chinese DLNM work is
urban --- but with six correlated pollutants and an interaction screen the multiple-testing
burden is not addressed, and the PM findings are the weakest of the set.}
{10.1080/09603123.2026.2718327}

\band{Mechanistic toxicology\hfill 1 record}{Violet}

\paperentry{Ayala et al.\ 2026}
{Contrasting neurotoxic pathways triggered by PM$_{10}$ and organic molecular markers, Catalonia}
{Environ Pollut}
{Fifty-four PM$_{10}$ filter samples collected through 2022 at a rural background site
(Bellver de Cerdanya), a suburban site (Manlleu) and a suburban-industrial site (Mollet
del Vall\`es), analysed by GC-MS for 30 organic tracers and tested on SH-SY5Y
neuroblastoma cells across six endpoints (viability, ROS, AChE activity, antioxidant and
xenobiotic response, p53). MCR-ALS on the joint chemical-biological matrix resolved four
components: winter biomass burning (levoglucosan, dehydroabietic acid, PAHs) driving
cytotoxicity, oxidative stress and xenobiotic response; a year-round traffic component; a
spring--summer SOA component producing \textbf{selective AChE inhibition without
cytotoxicity}; and a summer POA component. PLS models were significant for five of six
endpoints ($R^2_{CV}$ 0.53--0.75, best for ROS, p53 and cell death). The strength is that
exposure and effect were measured on the same extracts; the limits are one year, three
sites, one immortalised cell line and organic-solvent extracts that do not represent
inhaled particle dosimetry.}
{10.1016/j.envpol.2026.128961}

\band{Cardiovascular \& metabolic\hfill 1 record}{Teal}

\paperentry{Wittkopp et al.\ 2026}
{Cleaner Air for Lower Cardiometabolic Risk: HEPA filtration trial protocol in prediabetes}
{Am Heart J}
{Protocol for a double-blind, randomised, sham-controlled trial of one month of bedroom
HEPA filtration in adults with prediabetes in the New York City area (NCT05994937), with
continuous indoor PM$_{2.5}$ measured throughout. Primary outcomes are continuous glucose
monitoring metrics before and after; exploratory outcomes include insulin resistance,
serum biomarkers and transcriptomics. Secondary analyses will model measured PM$_{2.5}$
as a continuous independent variable, which is the analysis that will actually be
informative --- the intention-to-treat contrast in filtration trials is routinely diluted
by baseline indoor concentrations already being low. One month is far too short for any
glycated-haemoglobin endpoint, which is presumably why continuous glucose monitoring was
chosen instead. No results; logged to \texttt{trials.json}.}
{10.1016/j.ahj.2026.107573}

\band{Occupational \& indoor exposure\hfill 1 record}{Clay}

\paperentry{O'Donnell et al.\ 2026}
{Nicotine replacement therapy to reduce children's second-hand smoke exposure at home}
{medRxiv (preprint)}
{Pilot randomised controlled trial in Scotland testing whether free NRT plus
telephone-delivered support is a feasible route to a smoke-free home, with objective
indoor PM$_{2.5}$ as the exposure measure rather than self-reported smoking --- which is
the design decision that makes it relevant here. Feasibility, recruitment and retention
are the outcomes; the trial is not powered for exposure reduction. Notable as one of very
few interventions targeting the indoor source directly rather than filtering downstream
of it. Preprint, not peer reviewed; no numeric claim from it is used in this issue's
analytics.}
{10.64898/2026.08.13.26360348}

\band{Burden, policy \& mitigation\hfill 1 record}{Coral}

\paperentry{Wu \& Cui 2026}
{Aircraft CO, NO$_2$ and SO$_2$ emissions and public health at Chinese airports}
{J Glob Health}
{Landing-and-take-off cycle data from the Civil Aviation Administration of China,
concentrations at each airport from a Gaussian diffusion model, premature deaths from the
Global Exposure Mortality Model applied to populations near airports. Attributable
mortality was higher in men than women and concentrated in older age groups; CO accounted
for the most respiratory deaths; the peak burden year was 2019. \textbf{No particulate
metric is estimated at all}, which for an aviation source --- where ultrafine particle
number is the exposure of concern and is not proxied by CO --- is a substantive omission
rather than a scope choice. A Gaussian plume model at airport scale with no measurement
validation is the second limitation. Carried at tier C for the burden framing and as a
marker of where aviation-adjacent monitoring is not being done.}
{10.7189/jogh.16.04210}

\band{Other clinical endpoints\hfill 1 record}{Slate}

\paperentry{Negahi et al.\ 2026}
{Climate-sensitive exposures and cancer outcomes: systematic review and meta-analysis}
{Int J Biometeorol}
{Five databases searched to May 2026 for observational studies of wildfire, ambient
temperature, hurricane or multi-hazard disaster, precipitation anomaly, or ambient
PM$_{2.5}$ against cancer incidence, mortality, survival or recurrence; 25 estimates from
14 studies pooled within ten exposure-outcome groups, with robust variance estimation as
the primary method for groups containing dependent estimates. \textbf{Uncorrected pooled
1.07 (1.02--1.12); RVE 1.02 (0.93--1.11)} --- the conventional estimate was inflated by
correlated effects, with $I^2 = 97\%$. Wildfire and cancer mortality \textbf{1.11
(0.89--1.38)}; ambient PM$_{2.5}$ associated with poorer survival; an apparently
protective temperature-survival association vanished under RVE and is best read as
residual confounding. GRADE certainty was low in two groups and very low in the other
eight. Fourteen studies across five exposure types and four outcome classes is too thin a
base for group-specific inference, and the paper says so.}
{10.1007/s00484-026-03293-3}

\clearpage

% ==================================================================== TABLE
\section*{\sffamily\bfseries\color{Deep}Corpus register}
\vspace{-1.5mm}{\color{Amber}\rule{\linewidth}{1.1pt}}\vspace{2.5mm}

{\fontsize{6.7}{8.4}\selectfont
\setlength{\tabcolsep}{2.4pt}
\begin{longtable}{@{}>{\RaggedRight}p{25mm}>{\RaggedRight}p{41mm}>{\RaggedRight}p{27mm}%
>{\RaggedRight}p{23mm}>{\RaggedRight}p{18mm}>{\RaggedRight}p{15mm}@{}}
\toprule
\textsf{\textbf{First author}} & \textsf{\textbf{Focus}} & \textsf{\textbf{Journal}} &
\textsf{\textbf{Design}} & \textsf{\textbf{Metric}} & \textsf{\textbf{Region}}\\
\midrule
\endfirsthead
\toprule
\textsf{\textbf{First author}} & \textsf{\textbf{Focus}} & \textsf{\textbf{Journal}} &
\textsf{\textbf{Design}} & \textsf{\textbf{Metric}} & \textsf{\textbf{Region}}\\
\midrule
\endhead
\multicolumn{6}{@{}l@{}}{\cellcolor{Mist}\textsf{\textbf{\textcolor{Deep}{Neuro / mental health}}}}\\[0.6mm]
\href{https://doi.org/10.1001/jamanetworkopen.2026.25015}{Dales} & Wildfire-sourced PM$_{2.5}$ and ED visits for migraine and other primary headache & \textit{JAMA Netw Open} & Case-crossover & PM$_{2.5}$ (wildfire) & Canada\\
\href{https://doi.org/10.1093/ageing/afag215}{Glover} & Umbrella meta-review of urban environmental exposures and cognitive health & \textit{Age Ageing} & Umbrella review / meta-analysis & PM$_{2.5}$ & Global\\
\href{https://doi.org/10.1016/j.nsa.2026.107019}{Goossens} & In utero air pollution exposure and pre-/postnatal brain development & \textit{Neurosci Appl} & Narrative review & PM$_{2.5}$ / PM$_{10}$ / PM$_{\mathrm{coarse}}$ & Global\\
\href{https://doi.org/10.1016/j.isci.2026.116617}{Ma} & PM$_{2.5}$-induced neurotoxicity and depression: exposure risks and mechanisms & \textit{iScience} & Narrative review & PM$_{2.5}$ + components & Global\\
\href{https://doi.org/10.3760/cma.j.cn112150-20250804-00752}{Pan} & Maternal PM exposure and offspring neurodevelopment: effects and mechanisms & \textit{Zhonghua Yu Fang Yi Xue Za Zhi} & Narrative review & PM (bulk) & China\\
\addlinespace[0.9mm]
\multicolumn{6}{@{}l@{}}{\cellcolor{Mist}\textsf{\textbf{\textcolor{Teal}{Cardiovascular \& metabolic}}}}\\[0.6mm]
\href{https://doi.org/10.3760/cma.j.cn112338-20251120-00843}{Huang} & PM$_{2.5}$ and dyslipidaemia in rural Chinese children; dietary effect modification & \textit{Zhonghua Liu Xing Bing Xue Za Zhi} & Cross-sectional (multistage) & PM$_{2.5}$ & China\\
\href{https://doi.org/10.32394/pe/218966}{Kazmierska} & Air pollution and cardiovascular disease: review of current data & \textit{Przegl Epidemiol} & Narrative review & PM$_{2.5}$ / PM$_{10}$ & Global\\
\href{https://doi.org/10.1007/s00125-026-06814-2}{Xi} & Air pollution, built environment and temperature vs. prediabetes phenotypes & \textit{Diabetologia} & Prospective cohort & PM$_{2.5}$ & Germany\\
\href{https://doi.org/10.1161/JAHA.125.046917}{Yang} & Short-term PM$_{2.5}$ and hospitalisation/ED use in atrial fibrillation & \textit{J Am Heart Assoc} & Case-crossover & PM$_{2.5}$ (wildfire-influenced) & USA\\
\addlinespace[0.9mm]
\multicolumn{6}{@{}l@{}}{\cellcolor{Mist}\textsf{\textbf{\textcolor{Sky}{Reproductive \& developmental}}}}\\[0.6mm]
\href{https://doi.org/10.1016/j.jes.2025.10.043}{Chen} & Prenatal air pollution, fetal growth trajectory phenotypes and child neurodevelopment & \textit{J Environ Sci (China)} & Prospective cohort & PM$_{2.5}$ / PM$_{10}$ & China\\
\href{https://doi.org/10.1016/j.ecoenv.2026.120566}{Huang} & Joint prenatal bisphenol + air-pollutant mixture exposure and preschool IQ & \textit{Ecotoxicol Environ Saf} & Prospective cohort (mixtures) & PM$_{2.5}$ / PM$_{10}$ & China\\
\href{https://doi.org/10.1007/s40572-026-00552-8}{Peterson} & Ambient air pollution, PAHs and pubertal development: critical review & \textit{Curr Environ Health Rep} & Critical review & PM$_{2.5}$ components & Global\\
\href{https://doi.org/10.1016/j.envres.2026.125313}{Xie} & Long-term PM$_{2.5}$ and timing of pubertal development in boys and girls & \textit{Environ Res} & Prospective cohort & PM$_{2.5}$ & China\\
\addlinespace[0.9mm]
\multicolumn{6}{@{}l@{}}{\cellcolor{Mist}\textsf{\textbf{\textcolor{Sage}{Respiratory \& allergic}}}}\\[0.6mm]
\href{https://doi.org/10.1007/s00484-026-03289-z}{Bayar Muluk} & Climate change and otolaryngology: airway disease impacts and emission reduction & \textit{Int J Biometeorol} & Narrative review & PM$_{2.5}$ & Global\\
\href{https://doi.org/10.4103/mgr.MEDGASRES-D-25-00164}{Liao} & Short-term air pollution, respiratory symptoms and lung function in asthma & \textit{Med Gas Res} & Panel study & PM$_{2.5}$ / PM$_{10}$ & China\\
\href{https://doi.org/10.3390/epidemiologia7040097}{Sanchez de Prada} & Ambient PM and asthma-coded admissions in children 0-3 y, Spain & \textit{Epidemiologia} & Ecological panel & PM$_{2.5}$ / PM$_{10}$ & Spain\\
\href{https://doi.org/10.1016/j.jhazmat.2026.143076}{Yang} & Allergenic vs. toxic risk of airborne seed hairs compared with ambient PM$_{2.5}$ & \textit{J Hazard Mater} & Field measurement & PM$_{2.5}$ + bioaerosol & China\\
\addlinespace[0.9mm]
\multicolumn{6}{@{}l@{}}{\cellcolor{Mist}\textsf{\textbf{\textcolor{Violet}{Mechanistic toxicology}}}}\\[0.6mm]
\href{https://doi.org/10.1016/j.jes.2025.12.074}{Ding} & PM$_{2.5}$ promotes vascular calcification via miR-27b-5p/HES1 signalling & \textit{J Environ Sci (China)} & In vitro + organoid + murine & PM$_{2.5}$ & China\\
\href{https://doi.org/10.1007/s00360-026-01698-5}{Maraschi} & Osmoregulatory disruption and metal bioaccumulation from settleable atmospheric PM & \textit{J Comp Physiol B} & Controlled exposure (ecotox) & Settleable APM & Brazil\\
\href{https://doi.org/10.1016/j.jes.2025.11.018}{Wu} & Macrophage NLRP3 inflammasome in real-ambient PM-induced hepatic glucose dysregulation & \textit{J Environ Sci (China)} & Real-ambient murine + co-culture & Real-ambient PM & China\\
\addlinespace[0.9mm]
\multicolumn{6}{@{}l@{}}{\cellcolor{Mist}\textsf{\textbf{\textcolor{Clay}{Occupational \& indoor}}}}\\[0.6mm]
\href{https://doi.org/10.3389/fmedt.2026.1796828}{Chen} & Nested suction instrument to cut intraoperative surgical-smoke PM$_{2.5}$ & \textit{Front Med Technol} & Bench / device experiment & PM$_{2.5}$ (surgical smoke) & China\\
\href{https://doi.org/10.1093/annweh/wxag055}{Cole} & Determinants of respirable crystalline silica exposure in Australian tunnelling & \textit{Ann Work Expo Health} & Observational exposure study & Respirable dust / RCS & Australia\\
\href{https://doi.org/10.1016/j.jhazmat.2026.143081}{Ding} & Stage-resolved hair Zn-isotope dynamics as a firefighting combustion-exposure tracer & \textit{J Hazard Mater} & Repeated-measures biomarker & Combustion-derived PM & China\\
\href{https://doi.org/10.3390/jox16040119}{Ramsden} & Xenobiotic hazards in aircraft cabin air & \textit{J Xenobiot} & Review + risk economics & Ultrafine particles & Global\\
\addlinespace[0.9mm]
\multicolumn{6}{@{}l@{}}{\cellcolor{Mist}\textsf{\textbf{\textcolor{Amber}{Sensing, forecasting \& instrumentation}}}}\\[0.6mm]
\href{https://doi.org/10.1016/j.jenvman.2026.130515}{Shi} & Adaptive spatiotemporal graph transformer for multi-site multi-step PM$_{2.5}$ forecasting & \textit{J Environ Manage} & Deep-learning model & PM$_{2.5}$ & China / India\\
\href{https://doi.org/10.3390/bios16070380}{Tsuchiya} & Seasonal and indoor-outdoor characteristics of fluorescent aerosol particles in offices & \textit{Biosensors} & Real-time sensor campaign & Bioaerosol / coarse & Japan\\
\addlinespace[0.9mm]
\multicolumn{6}{@{}l@{}}{\cellcolor{Mist}\textsf{\textbf{\textcolor{Amber}{Exposure assessment \& modelling}}}}\\[0.6mm]
\href{https://doi.org/10.1029/2025GH001621}{Solorzano-Araque} & Multiscale exposure modelling: LOTOS-EUROS + mobility surveys + particle toxicity & \textit{GeoHealth} & CTM + mobility integration & PM$_{2.5}$ / PM$_{10}$ & Colombia\\
\addlinespace[0.9mm]
\multicolumn{6}{@{}l@{}}{\cellcolor{Mist}\textsf{\textbf{\textcolor{Coral}{Burden, policy \& mitigation}}}}\\[0.6mm]
\href{https://doi.org/10.1002/hsr2.72886}{Karamzad} & Ischaemic heart disease burden in MENA 1990-2021; PM attributable fraction 28.2\% & \textit{Health Sci Rep} & GBD secondary analysis & Ambient PM & MENA\\
\href{https://doi.org/10.1371/journal.pone.0347989}{Rees} & HEPA filtration to reduce respiratory infection in care homes & \textit{PLoS One} & RCT process evaluation & Filtered PM & UK\\
\href{https://doi.org/10.1016/j.jes.2026.01.009}{Shan} & Policy-driven emission reductions and post-pandemic rebound, high-resolution inventories & \textit{J Environ Sci (China)} & Emission inventory & PM emissions & China\\
\addlinespace[0.9mm]
\multicolumn{6}{@{}l@{}}{\cellcolor{Mist}\textsf{\textbf{\textcolor{Slate}{Other clinical endpoints}}}}\\[0.6mm]
\href{https://doi.org/10.3390/dj14070440}{Lupita} & PM$_{2.5}$, preventive-care capacity and caries/edentulism burden across the EU & \textit{Dent J} & Ecological panel & PM$_{2.5}$ & EU-27\\
\href{https://doi.org/10.1016/j.jdin.2026.06.008}{Park} & Long-term PM exposure and risk of androgenetic alopecia & \textit{JAAD Int} & Nationwide cohort & PM$_{2.5}$ / PM$_{10}$ & South Korea\\
\href{https://doi.org/10.3171/2026.3.PEDS2615}{Sun} & Environmental exposure and area deprivation vs. paediatric hydrocephalus outcomes & \textit{J Neurosurg Pediatr} & Registry cohort + geospatial & PM$_{2.5}$ (satellite-derived) & USA\\
\addlinespace[0.9mm]
\bottomrule
\end{longtable}}


\vspace{4mm}
\begin{keybox}
{\sffamily\bfseries\footnotesize\color{Deep}Provenance and method}\\[1.2mm]
{\fontsize{7.2}{8.8}\selectfont
\textbf{Window.} PubMed \texttt{[EDAT]} \textbf{2026-08-14 to 2026-08-14}, contiguous with
the 13 August issue; no gap. Every record was checked against
\texttt{state/seen.json} (pmid $\rightarrow$ lowercased doi $\rightarrow$ title signature)
before summarising; no cross-source duplicate arose this issue.\\[0.8mm]
\textbf{Per-leg yield.} PubMed connector, two separate queries on the health and
instrumentation axes: 8 and 5 hits, 13 distinct, \textbf{9 carried}. Local harvester:
\texttt{pubmed\_health} 15 and \texttt{pubmed\_sensing} 1, of which five records were
outside the connector's two queries and \textbf{2 of those were carried} (Zhai \& Meng;
Wu \& Cui) --- the fourth consecutive issue in which the local leg has earned its
keep; Europe PMC 3 (\textbf{2 carried}, both preprints); Crossref by ISSN 7
(\textbf{1 carried} --- three were cloud-microphysics papers from \textit{Atmos. Chem.
Phys.}, one was an \textit{Atmos.\ Pollut.\ Res.} editorial board listing, one an
Elsevier metadata-only VOC record). arXiv returned 146\,kB with no entry on or after
10 August. \textbf{OpenAlex HTTP 429, 12th consecutive} --- the leg remains broken and is
not worth further effort. \textbf{Consensus} was re-queried on the calibration and drift
axis and returned three highly relevant papers, \textbf{all out of window} (Crossref
\texttt{created} 2026-02-19, 2026-03-26 and 2026-06-15); this is the seventh consecutive
issue with zero Consensus carry, and the pattern is now clear enough to record: Consensus
ranks by relevance, not recency, and cannot serve a one-day entry window.
ClinicalTrials.gov: no PM or air-pollution registration created or updated inside the
window; \texttt{trials.json} unchanged, though NCT05994937 now has a published protocol
(Wittkopp et al., this issue).\\[0.8mm]
\textbf{Rejections (17).} Fourteen in-harvest: one published erratum; one \textit{JAMA}
Medical News item reporting a study already covered; three non-PM toxicology papers
(chlorpyrifos, ingested polypropylene microplastics + DEHP, ethanol neuroinflammation); a
greenhouse-gas inventory for source-separated organics; a fly-ash/farmyard-manure
agronomy study; a scoping review whose literature search closed in April 2021; two
cloud-microphysics papers and one ice-nucleating-protein paper; a Permian Basin
hydrocarbon source-apportionment study with no PM metric and no abstract released; a
journal editorial board listing; and an editorial commentary on emerging POPs. Three
further records were in scope but outside the entry window (the Consensus hits above).
All are logged with reasons to \texttt{state/rejected.jsonl}.\\[0.8mm]
\textbf{Two records are preprints} (Research Square, medRxiv) and are labelled as such in
the register. Neither supports a numeric claim in the analytics.\\[0.8mm]
\textbf{Labels.} Subtopic, design, particle-metric and geography labels were drawn only
from vocabularies already mapped in \texttt{plots.py} and \texttt{mktable.py}; no new leaf
label was written this issue. Assignment is single-label by dominant contribution and is a
judgement, not a controlled vocabulary. \texttt{check\_dois.py} returned 0 fail, 0 warn
across all 14 records.\\[0.8mm]
Content derived from PubMed, Europe PMC and Crossref metadata; abstracts remain the
copyright of their respective publishers. All DOI links resolve to the publisher of
record.\par}
\end{keybox}

\end{document}
